% Nouveau EGA V (cote 25): archive cote 25, site page 114 (PDF page 115).
% typescript. Draft compared: none.
% AI-transcribed by Claude Opus 5.5 (claude-opus-5-5, adaptive thinking, effort high), Message Batches API, batch ega25b (msgbatch_01A7n5YuXubUPNK1zsepi8RG), 2026-11-13; not yet checked by eye.
\scanpage{25/114}
\mspage{17}

On pourra noter aussi en remarque que lorsque X et T sont tous les
deux plats loc de présentation finie sur S, à fibres $(S_1)$ et géom irréd.,
alors la symétrie $D \rightsquigarrow {}^{t}D$ donne une correspondance biunivoque entre
familles de diviseurs de X/S paramétrées par T qui sont sans points fixes
et familles de diviseurs de T/S paramétrées par X qui sont sans points
fixes. Si dans cet énoncé on désire se débarasser des hypothèses spécifi-%
ques faites sur les fibres de X/S et T/S , il convient de remplacer
les ``points fixes'' par ``points fixes au sens large'' , en entendant par
point fixe au\xout{x} sens large de D un $x \hand{\in} X$ tel que D ne soit pas diviseur
\deleted{de Cartier} relatif \xout{en} \inserted[typed]{par rapport à X en tous} les points de $\mathrm{pr}_1^{-1}(x)$ . Si W est l'ouvert des
points de $X\times_S T$ en lesquels D \xout{n'}est \xout{pas} diviseur \deleted{de Cartier} relatif par
rapport à X, alors l'ensemble des points fixes au s ens large de D est
\xout{\unsure{exemple}} égal à $\mathrm{pr}_1(X\times_S T - W)$; lorsque $T \to S$ est propre, c'est donc
une partie fermée Z' de X, \xout{et on obtient} En tous cas, on obtient une
famille de diviseurs de T/S paramétrée par $X-\overline{Z'}$ . L'hypothèse que les
fibres de T/S soient $(S_1)$ et géom irréductibles \xout{\unsure{assure}} \inserted[typed]{sert} précisément à
assurer que $Z=Z'$ (point fixes \unsure{$=$} points fixes au sens large). Géométri-%
quement, supposant pour simplifier que S=Spec(k), k corps algébriquement
clos (ce qui est loisible, \inserted[typed]{pour T/S plat loc de prés finie,} via changement de base), dire que $x \hand{\in} X(k)$
est point fixe (resp. point fixe au sens large) signifie que $x \hand{\in} \mathrm{supp}\, D_t$
pour tout $t \hand{\in} T(k)$ (resp. qu'il existe un cycle premier T' associé à T
tel que $x \hand{\in} D_t$ pour tout $t \hand{\in} T'(k)$).

\deleted{Le petit sorite précédent n'a guère rien à voir dansle présent para-%
graphe, et je suggérerai d'en faire plutôt un nº au par. 19 à propos des
intersections complètes relatives, sous le titre: diviseurs de Cartier
relatifs, et familles \xout{algébriques} de diviseurs.} \margin{$\|$} Un oubli: \xout{X-Z est}
La formation de l'ensemble des points fixes Z est compatible avec

\mspage{114}
